% Figure: bias squared, variance, and total MSE as a function of the
% ridge regularisation parameter lambda. Stylised curves showing the
% canonical U-shape of the MSE.
\begin{figure}[htbp]
\centering
\begin{tikzpicture}
\begin{axis}[
    width=11cm, height=6cm,
    xmin=0.05, xmax=100,
    ymin=0, ymax=2.5,
    axis lines=left,
    xmode=log,
    log basis x=10,
    xlabel={Regularisation parameter $\lambda$ (log scale)},
    ylabel={Error component},
    samples=200,
    domain=0.05:100,
    no markers,
    every axis plot/.append style={very thick},
    legend style={font=\small, draw=none, fill=none, at={(0.02, 0.98)}, anchor=north west},
    ytick=\empty,
]
\addplot[engblue]   {2/(1+x)};
\addlegendentry{Variance $V(\lambda)$}

\addplot[engred]    {0.05*x*x/(1+x)/(1+x)*2};
\addlegendentry{Bias$^2$ $B^2(\lambda)$}

\addplot[engblack, ultra thick]
                    {2/(1+x) + 0.05*x*x/(1+x)/(1+x)*2};
\addlegendentry{MSE $= V + B^2$}

\addplot[enggray, dotted] coordinates {(4, 0) (4, 2.5)};
\node[anchor=south, font=\small, engblack] at (axis cs:4, 1.1) {$\lambda^{\star}$};
\end{axis}
\end{tikzpicture}
\caption[Bias-variance trade-off under ridge regularisation]{The
bias-variance decomposition for the ridge regression estimator as
the penalty $\lambda$ varies. Variance decreases monotonically with
$\lambda$; squared bias increases monotonically; total mean squared
error has an interior minimum at $\lambda^{\star}$. The shape is
the canonical justification for biased estimation. Ridge, lasso, and
Bayesian shrinkage all sit on a curve of this kind. The choice of
$\lambda$ in practice is made by cross-validation, not by visual
inspection of a plot the engineer drew on training data.}
\label{fig:vol02:ch14:bias-variance}
\end{figure}
